\section{Introdução}

Em plataformas de transporte, o deslocamento do veículo até o embarque acrescenta quilômetros sem passageiro à operação. Neste estudo, \textit{km vazio} é apenas a distância estimada de aproximação aos pedidos despachados; não inclui reposicionamento nem retorno. Como cada escolha altera a disponibilidade e a posição futura da frota, atribuir isoladamente o veículo elegível mais próximo pode afetar os atendimentos seguintes. Investiga-se a decisão de formar pares entre pedidos pendentes e veículos livres no lote corrente de 60 segundos, priorizando o número de pares e, depois, a soma das distâncias de aproximação, sem prever demanda futura.

A pergunta de pesquisa é: \textbf{em simulações com os mesmos pedidos, frota, jornadas e método de distância, a atribuição conjunta por lotes de 60 segundos reduz o km vazio sem prejudicar o atendimento, em comparação com o despacho imediato ao veículo elegível mais próximo e com o despacho sequencial em lote?}

A hipótese fixada antes das execuções finais é que a atribuição conjunta pelo algoritmo Húngaro \cite{kuhn1955, munkres1957} reduz o km vazio nas duas comparações. A avaliação dessa redução exige também atendimento preservado: ausência de violação de compromisso causada pela política, pelo menos tantos pedidos atendidos quanto na referência e percentil 90 da espera não superior ao da referência. Trata-se da expectativa inicial, a ser confrontada com resultados favoráveis e desfavoráveis.

O objetivo geral deste trabalho é avaliar, por simulação, como a atribuição conjunta em lotes afeta o km vazio e o atendimento frente ao despacho imediato e ao despacho sequencial em lote. A comparação com este último examina a decisão conjunta sob a mesma espera pelo lote; a comparação com o despacho imediato abrange tanto a espera quanto a escolha dos pares.

Para alcançar esse objetivo, busca-se formular a atribuição por lote, implementar e verificar as políticas, estruturar avaliações com entradas comuns, comparar distância, atendimento e espera e examinar os limites do simulador frente aos registros da operação. Os períodos históricos de setembro e outubro não participaram da escolha do método; cenários com demanda e frota sintéticas fornecem evidência exploratória separada.

O problema é relevante porque uma redução de aproximação pode coexistir com perda de serviço, e a escolha ótima em um lote não garante melhor resultado ao longo do dia. A contribuição acadêmica é distinguir o efeito de esperar pelo lote do efeito de atribuir conjuntamente, com critérios de atendimento explícitos e verificações do modelo. A relevância prática está em informar quais benefícios e limites exigiriam validação antes de qualquer uso operacional; combustível, emissões e economia financeira não foram medidos.

A investigação usa simulação de eventos discretos com pedidos, jornadas e posições históricas de uma empresa de transporte por aplicativo na Grande Vitória (ES), além de cenários sintéticos sobre a rede viária da região. Comparam-se despacho imediato, lote sequencial e atribuição conjunta sob as mesmas entradas de cada avaliação; o guloso global integra apenas análise posterior. As medidas principais são km vazio estimado e atendimento, com espera como controle. Os resultados descrevem diferenças no simulador, não uma economia observada na operação.
